% ===================================================================
% CHAPTER 4: RELATED WORK
% ===================================================================
\chapter{Поврзани истражувања и литература}
\label{chap:related}

\section{Историски контекст и еволуција на федерираното откривање упади}

Системите за откривање упади засновани на машинско учење поминаа низ три фази. Првата беше
целосно централизирана: сообраќајот од сите точки на набљудување се собира на едно место и
таму се обучува моделот. Пристапот дава највисока точност, но е неприменлив во опкружувања
со Интернет на нештата --- преносот на суров мрежен сообраќај е скап, ја нарушува
приватноста и создава единствена точка на дефект.

Втората фаза го воведе федерираното учење како одговор \cite{mcmahan2017fedavg}. Наместо
податоци, се пренесуваат ажурирања на параметри. За откривање упади тоа е особено
привлечно: обележјата на мрежните текови се чувствителни, а самата задача --- бинарна
класификација \emph{напад} наспроти \emph{бенигно} --- е доволно едноставна за да се изведува и на
скромен хардвер.

Практичната пречка во таа фаза е \textbf{хетерогеноста}. Во реална поставка секој јазол гледа
сообраќај од различен извор, па локалните множества не се независни и идентично
распределени. FedProx \cite{li2020fedprox} го решава тоа со проксимален член што ја
ограничува оддалеченоста на локалното решение од глобалното. Овој труд ја наследува истата
особеност по конструкција: поделбата се врши на ниво на \textbf{список од датотеки}, а не на
редови, така што секој јазол добива дисјунктно множество целосни снимки на сообраќај и
партициите се природно non-IID по извор.

Третата фаза, во која се наоѓа ова истражување, го поместува тежиштето од \emph{обука} кон
\textbf{одржување}: што се случува со федерираниот модел откако ќе се утврди дека дел од
податоците не смеат да бидат во него.

\section{Преглед на методологии за федерирано заборавање}

Машинското заборавање има јасен почеток во работата на Cao и Yang \cite{cao2015unlearning},
кои покажаа дека за широк класа алгоритми обуката може да се преформулира во збирна форма
што дозволува ефикасно отстранување на поединечен придонес. Општата рамка за егзактно
заборавање на произволни модели ја даде SISA \cite{bourtoule2021sisa}, со шардирање и
сегментирање што однапред го ограничува досегот на секој примерок. Прегледите на областа
\cite{nguyen2022musurvey} ја систематизираат поделбата на егзактни и приближни постапки.

Пренесувањето во федерираната поставка е предмет на низа поскорешни трудови, кои прегледите
\cite{liu2024fusurvey, romandini2024fusurvey} ги групираат по тоа \textbf{каква состојба се чува
и кај кого}:

\begin{itemize}
  \item \textbf{Реконструкција од историски ажурирања.} FedEraser \cite{liu2021federaser} ги
  калибрира зачуваните ажурирања на преостанатите клиенти за да ја реконструира траекторијата
  на обуката без придонесот на отстранетиот клиент, што е значително побрзо од целосна
  повторна обука. FedRecover \cite{cao2023fedrecover} го применува истиот принцип на
  обновување по напад со труење, користејќи ги историските информации за да го поправи
  глобалниот модел со мал број пресметки кај клиентите.
  \item \textbf{Шардирање на ниво на клиенти.} Директното пресликување на SISA во федерираниот
  простор ги групира клиентите во дисјунктни кластери и обучува независен модел по кластер;
  при барање за заборавање се преобучува само засегнатиот кластер. Lin и соработниците
  \cite{lin2024codedsharding} покажуваат дека комбинацијата на изолирано шардирање со
  кодирани параметри го намалува времето на преобука и складирањето во однос на FedEraser и
  на целосната повторна обука.
\end{itemize}

Заедничкото за сите нив е една силна претпоставка: \textbf{серверот учествува} во заборавањето ---
или чува историја на ажурирања, или ја менува шемата на агрегација, или ја наметнува
организацијата на клиентите во кластери. Таа претпоставка е разумна во истражувачка
поставка, но е сериозна пречка при вградување во постоечка федерација што веќе работи со
стандарден FedAvg.

Второ заедничко својство е дека сите наведени трудови ја изразуваат цената во
\textbf{пресметковни} единици --- рунди, епохи, комуникациски трошок --- и ја мерат во симулација.

\section{Труење со податоци и робусна агрегација}

Отвореноста на федерираното учење е истовремено и негова ранливост: придонесот на
учесникот по правило се прифаќа без проверка. Tolpegin и соработниците
\cite{tolpegin2020poisoning} покажуваат дека насочено труење со превртување на етикети, и
со мал удел злонамерни учесници, предизвикува значителен пад на одѕивот за целната класа.
Во контекст на откривање упади тоа е особено опасно: превртувањето \emph{напад $\to$ бенигно} е
токму дефиницијата на задна врата во детекторот.

Одбраната по правило се бара во робусната агрегација --- Krum \cite{blanchard2017krum},
скратената средна вредност и медијаната по координата \cite{yin2018trimmedmean}. Fang и
соработниците \cite{fang2020localpoisoning} потоа покажуваат дека тие правила може да се
победат со напади скроени токму за нив, што ја релативизира нивната заштитна улога.

За овој труд е важна една концептуална разлика, која литературата за робусност често ја
замаглува: \textbf{робусната агрегација го ублажува влијанието, но не го отстранува податокот}.
Krum и скратената средна вредност ја штитат точноста на моделот; тие не даваат основа да се
тврди пред регулатор дека определени податоци повеќе не се во него. Тоа е предмет на
заборавањето. Робусните стратегии беа испитани во претходната фаза на овој проект (FedAvg,
FedProx, Krum и скратена средна вредност врз истото податочно множество) и служат како
контекст, но не се алтернатива на постапката што се мери овде.

\section{Потребата за хардверски-свесна евалуација}

Прегледот погоре открива конзистентна \textbf{празнина}.

Литературата за федерирано заборавање ја мери цената на отстранувањето во епохи, рунди на
комуникација и операции со подвижна запирка, и речиси исклучиво во симулација на сервер или
на работна станица. Литературата за федерирано откривање упади на уреди од класата
Raspberry Pi, од друга страна, известува посредни показатели --- искористеност на процесорот
и меморијата, температура, време на обука по рунда, доцнење при заклучување --- но ја мери
\emph{обуката}, не \emph{заборавањето}, и ретко ја мери потрошувачката со надворешен инструмент.

Со други зборови: постојат трудови што го мерат заборавањето, и постојат трудови што мерат
физички ресурси на работ, но речиси нема трудови што ги прават двете истовремено.
Последицата е практична. Проектант што треба да одлучи дали да ја плати непрекинатата цена
на SISA --- запишување контролни точки во секоја рунда, на SD-картичка со ограничена
издржливост --- за да добие побрзо обновување подоцна, во постојната литература нема на што
да се потпре. Прашањата што остануваат без емпириски одговор се конкретни:

\begin{enumerate}
  \item Колку \textbf{ват-часови} навистина се штедат, наспроти теоретското намалување на бројот на епохи?
  \item Дали намаленото траење го спречува \textbf{термичкото придушување}, кое на уред без активно ладење може да ја обезличи секоја пресметка на хартија?
  \item Дали цената навистина се преместува кон \textbf{влезно-излезни операции}, и при која големина на моделот тоа станува доминантно?
  \item Дали обновениот модел останува \textbf{употреблив} --- мерено со показатели што издржуваат изразена класна нерамнотежа, а не со одѕив и F1?
\end{enumerate}

Овој труд ја пополнува таа празнина: ги спроведува двата пристапа на \textbf{физички} јазол на
работ, ги мери времето, енергијата, топлината и влезно-излезните операции со надворешни
инструменти, и ги подложува на претходно регистриран спарен експериментален дизајн со
непараметриска статистика. Придонесот не е нов алгоритам за заборавање, туку
\textbf{емпириска основа за избор} меѓу двата постоечки, изразена во величини што
проектантот на системи навистина ги плаќа.
